\documentclass{article}

\usepackage{booktabs}
\usepackage{tabularx}
\usepackage{longtable}
\usepackage{array}
\usepackage{pdflscape}
\usepackage{eso-pic}
\usepackage[margin=1in]{geometry}
\usepackage{hyperref}

\hypersetup{
    colorlinks=true,       % false: boxed links; true: colored links
    linkcolor=red,          % color of internal links (change box color with linkbordercolor)
    citecolor=green,        % color of links to bibliography
    filecolor=magenta,      % color of file links
    urlcolor=cyan           % color of external links
}

\title{Hazard Analysis\\\progname}

\author{\authname}

\date{}

\input{../Comments.text}
\input{../Common.text}

% Ragged-right fixed-width column for the FMEA table
\newcolumntype{P}[1]{>{\raggedright\arraybackslash}p{#1}}

\begin{document}

\maketitle
\thispagestyle{empty}

~\newpage

\pagenumbering{roman}

\begin{table}[hp]
\caption{Revision History} \label{TblRevisionHistory}
\begin{tabularx}{\textwidth}{llX}
\toprule
\textbf{Date} & \textbf{Developer(s)} & \textbf{Change}\\
\midrule
October 18, 2026 & Rhea Gokhale & Initial Revision 0 draft\\
\bottomrule
\end{tabularx}
\end{table}

~\newpage

\tableofcontents

~\newpage

\pagenumbering{arabic}

\section{Introduction}

This document is the hazard analysis for \progname{}, a language learning product based on songs. It is made up of a mobile learner application and a song map editor that runs in a web browser. Learners study synchronized, annotated lyrics on
their phone and send selected vocabulary to Anki. Creators use the editor to
prepare the song maps that drive each lesson.

\progname{} is not a system on which safety critically depends: it
controls no machinery and cannot directly injure anyone. It is, however, used by
the general public, often with headphones and on the move, and it writes into
the learner's own Anki collection. Its hazards are therefore mostly about
hearing and attention, incorrect learning content, loss or corruption of user
data, privacy, legal exposure from song content, and the learner's wellbeing.

\subsection*{Definition of Hazard}

Following Nancy Leveson's work, a \textbf{hazard} is a property or condition of
\progname{} that, together with a condition in its environment, can lead to a
\textbf{loss}. A \textbf{loss} is any outcome that a stakeholder would find
unacceptable. The losses considered in this document are listed in
Section~\ref{SecScope}.

\subsection*{Changes to the Template}

The template has been followed with two additions. Section~\ref{SecFMEA} adds a
\emph{Hazards Out of Scope} subsection, and the FMEA table rates each failure
mode for severity, occurrence, and detection and gives its risk priority
number (RPN). The RPN is used to prioritize the requirements in the Roadmap
(Section~\ref{SecRoadmap}).

\section{Scope and Purpose of Hazard Analysis} \label{SecScope}

The purpose of this analysis is to identify the ways the components of
\progname{} can fail or be misused, the effects and causes of those failures,
and the requirements or other mitigations that reduce each risk to an
acceptable level. Requirements discovered here are listed in
Section~\ref{SecRequirements} and are included in the Software Requirements
Specification (SRS).

The analysis covers both failures (the system does something wrong) and hazards
that exist when the system works as designed (for example, a correctly
synchronized lesson that distracts a learner who is walking in traffic).

The losses that the hazards in this document could lead to are:

\begin{itemize}
  \item \textbf{L1 Physical harm:} hearing discomfort or damage from
        unexpected audio, injury from distraction while moving, or discomfort
        for photosensitive users from flashing highlights.
  \item \textbf{L2 Incorrect learning:} the learner learns a wrong meaning,
        reading, pronunciation, or word boundary, or unknowingly learns an
        offensive term, and carries that error into real conversations and
        into their Anki reviews.
  \item \textbf{L3 Loss or corruption of user data:} notes in the learner's
        Anki collection are duplicated, malformed, or overwritten, selected
        vocabulary is silently lost, or a creator loses timing and annotation
        work.
  \item \textbf{L4 Privacy and security:} personal data leaves the device
        without consent, or a malicious song map or misconfigured integration
        allows an attacker to run code or change the learner's data.
  \item \textbf{L5 Legal:} distributing lyrics or recordings without the
        right to do so.
  \item \textbf{L6 Wellbeing:} the product causes stress, discouragement, or
        compulsive use, for example through manipulative engagement
        mechanics.
\end{itemize}

\section{System Boundaries and Components}

The system under analysis consists of the following components, which are
under the team's control:

\begin{enumerate}
  \item \textbf{Lesson Player} (learner app): plays the song's audio and
        highlights the current lyric line and word using the song map's
        timestamps.
  \item \textbf{Lyric Display and Language Support} (learner app): renders
        lyrics in their original script, the stored word segmentation, and the
        stored transliteration.
  \item \textbf{Annotation and Lesson Content}: translations, contextual
        annotations, and content metadata contained in each song map and
        shown to the learner.
  \item \textbf{Song Map Format and Validator}: the shared JSON Schema and the
        validation code used by both the editor (on export) and the learner
        app (on import).
  \item \textbf{Song Map Editor} (web): the creator's tool for aligning timing,
        entering translations and annotations, verifying segmentation and
        transliteration, and exporting song maps.
  \item \textbf{Anki Handoff}: the learner app's module that sends a selected
        vocabulary item to Anki through the AnkiDroid API, AnkiConnect, or an
        export file (the method is decided by the Proof of Concept, risk R3 in
        the Development Plan).
  \item \textbf{Local Data Storage} (learner app): stored lessons, display
        preferences, and the queue of vocabulary waiting to be sent to Anki.
\end{enumerate}

The following elements are outside the team's control but are still critical
to the system, so they are included as sources of failure where relevant:

\begin{itemize}
  \item the learner's phone and operating system (iOS or Android), including
        system volume, audio routing, and headphones,
  \item the creator's computer and web browser,
  \item Anki, AnkiDroid, and the AnkiConnect extension, and the learner's Anki
        collection,
  \item the audio and lyric content itself and the rights attached to it.
\end{itemize}

The initial release has no user accounts, no backend server, and no cloud
sync, as stated in the Problem Statement and Development Plan. Hazards specific
to those features are therefore not analyzed. They must be added if those
features are introduced.

\section{Critical Assumptions}

The team has kept assumptions to a minimum. In particular, song maps and their
creators are \emph{not} assumed to be trustworthy or free of errors, and learners
are \emph{not} assumed to use the app only while stationary. Both are treated
as hazards in Section~\ref{SecFMEA}.

\begin{enumerate}
  \item[CA1:] Software from other vendors (the phone operating system, the web
        browser, Anki, AnkiDroid, and AnkiConnect) behaves as described in its
        public documentation. Failures caused by defects inside that software
        are out of scope, although the app must still detect and report a
        failed handoff (H6.4).
  \item[CA2:] The learner's phone meets the minimum operating system version
        supported by the Expo SDK release the app is built with.
\end{enumerate}

\section{Failure Mode and Effect Analysis} \label{SecFMEA}

A Failure Modes and Effects Analysis (FMEA) was used to identify and analyze
hazards component by component. Each row lists a failure mode, its effects
(with the related loss from Section~\ref{SecScope}), its causes, how the
failure is detected, the recommended action, its ratings for severity,
occurrence, and detection with the resulting risk priority number, and the
safety or security requirements (Section~\ref{SecRequirements}) that mitigate
it. Each
row has a reference number (\emph{Ref.}) so that requirements can trace back
to it.

\subsection{Hazards Out of Scope}

\begin{itemize}
  \item Hardware failure of the learner's phone, headphones, or the creator's
        computer.
  \item Defects inside Anki, AnkiDroid, AnkiConnect, the operating system, or
        the browser (see CA1).
  \item The overall volume the learner chooses with their device's controls.
        The app controls only when playback starts and stops (H1.2). Limits on listening over long periods are managed by the operating system's hearing
        protection features.
  \item Hazards of streaming service integration, a production music
        catalogue, user accounts, and cloud sync, none of which are in the
        initial release.
  \item A user deliberately modifying the installed app or the files on their
        own device.
\end{itemize}

\subsection{Rating Scales}

Each failure mode is rated on three scales from 1 to 10. The ratings describe
the failure mode \emph{before} the recommended action is taken, so they show
how much each action matters.

\begin{table}[h!]
\caption{Rating scales} \label{TblScales}
\begin{tabularx}{\textwidth}{lXXX}
\toprule
\textbf{Rating} & \textbf{Severity (S)} & \textbf{Occurrence (O)} &
\textbf{Detection (D)}\\
\midrule
9 to 10 & Physical harm (L1), legal exposure (L5), or a privacy or security
breach (L4) & Almost certain in normal use & Cannot be detected before it
reaches the user\\
7 to 8 & Permanent loss or damage of user data (L3), or a serious privacy or
security weakness & Frequent & Detected only by chance\\
5 to 6 & The learner learns incorrect content (L2), or a goal from the
Problem Statement fails & Occasional & Detected by careful review or
targeted testing\\
3 to 4 & User data is lost but can be recovered, or the user is clearly
inconvenienced & Rare & Usually detected by routine tests\\
1 to 2 & Minor inconvenience with no lasting effect & Very unlikely & Almost
certainly detected by automated checks\\
\bottomrule
\end{tabularx}
\end{table}

The risk priority number is $\mathit{RPN} = S \times O \times D$, which ranges
from 1 to 1000. A failure mode is \textbf{high priority} if
$\mathit{RPN} \ge 150$ or $S \ge 9$. The second condition keeps a failure
mode that could cause physical harm or legal exposure high priority even when
it is unlikely. The ratings are estimates by the team and will be reviewed
after the Proof of Concept.

\subsection{FMEA Table}

The FMEA table is given on the following pages. Table~\ref{TblPriority}, after
it, lists the high priority failure modes in order of RPN.

\begin{landscape}
% Landscape pages are rotated for viewing, so the normal footer would appear
% on the side. Print the page number at the physical right edge instead,
% which is the bottom of the page once it is rotated.
\pagestyle{empty}
\AddToShipoutPictureFG{%
  \AtPageLowerLeft{%
    \put(\LenToUnit{\paperwidth-0.6in},\LenToUnit{0.5\paperheight}){%
      \rotatebox{90}{\makebox[0pt]{\normalsize\thepage}}}}}
\footnotesize
\setlength{\tabcolsep}{4pt}
\begin{longtable}{P{0.04\linewidth}P{0.10\linewidth}P{0.12\linewidth}P{0.11\linewidth}P{0.10\linewidth}P{0.19\linewidth}P{0.022\linewidth}P{0.022\linewidth}P{0.022\linewidth}P{0.03\linewidth}P{0.06\linewidth}}
\caption{FMEA Worksheet} \label{TblFMEA}\\
\toprule
\textbf{Ref.} & \textbf{Failure Mode} & \textbf{Effects of Failure} &
\textbf{Causes of Failure} & \textbf{Detection} & \textbf{Recommended Action} &
\textbf{S} & \textbf{O} & \textbf{D} & \textbf{RPN} & \textbf{SR}\\
\midrule
\endfirsthead
\multicolumn{11}{l}{\tablename\ \thetable{}, continued from previous page}\\
\toprule
\textbf{Ref.} & \textbf{Failure Mode} & \textbf{Effects of Failure} &
\textbf{Causes of Failure} & \textbf{Detection} & \textbf{Recommended Action} &
\textbf{S} & \textbf{O} & \textbf{D} & \textbf{RPN} & \textbf{SR}\\
\midrule
\endhead
\midrule
\multicolumn{11}{r}{Continued on next page}\\
\endfoot
\bottomrule
\endlastfoot
%
% ---------------------------------------------------------------- H1
\multicolumn{11}{l}{\textbf{H1: Lesson Player}}\\
\midrule
H1.1 & Lyric highlighting drifts from the audio by more than 0.5\,s. &
Learner connects the wrong words to the sounds they hear and may learn the wrong
pronunciation (L2). The synchronized playback goal fails. &
a.~Audio library reports position too coarsely or late.\newline
b.~App returns from the background, or the user seeks or pauses.\newline
c.~Timestamps in the map are wrong (see H5.3). &
Position logging in the acceptance suite of two songs. Developer debug overlay
showing drift. &
Drive highlighting from the player's reported position rather than a separate
timer, and realign immediately after every seek, pause, resume, and return from
the background. If the PoC (risk R1) shows positions are too coarse, switch to
a native audio module. &
6 & 7 & 3 & 126 & IN1, IN2\\
\midrule
H1.2 & Audio starts unexpectedly or at a sudden high level. &
Hearing discomfort or damage with headphones (L1). Audio plays aloud in a public
place after headphones are unplugged. &
a.~Lesson autoplays when opened.\newline
b.~Playback continues on the phone speaker when headphones disconnect.\newline
c.~Tracks are mastered at very different loudness. &
Manual test on iOS and Android: open lesson, unplug or disconnect headphones
during playback. &
Never start playback without a learner action. Pause playback when the audio
output device is disconnected. Fade audio in when playback starts. The editor
reports the loudness of each track so creators can normalize it before export.
&
8 & 5 & 4 & 160 & HS1, HS2\\
\midrule
H1.3 & Audio file referenced by a map is missing, corrupt, or in an unsupported
format. &
Lesson cannot be played. The app may crash (L3 if queued vocabulary is lost). &
a.~File not included in the map package.\newline
b.~Codec not supported on the learner's platform.\newline
c.~File damaged during transfer. &
Check on import that the file exists, decodes, and has a duration consistent
with the map. &
Reject the lesson on import with a message naming the problem, instead of
failing at playback time. Restrict maps to audio formats supported on both
iOS and Android. &
4 & 4 & 2 & 32 & IN11\\
\midrule
H1.4 & Learner watches the screen while walking, cycling, or driving. &
Physical injury to the learner or others from distraction (L1). &
a.~Intended use includes short study moments on the move.\newline
b.~Following the lesson requires reading highlighted text. &
Cannot be reliably detected by the app. &
Provide a listening mode that can be followed without looking at the
screen (larger controls, no required taps). Show a notice on first launch only, advising against using the screen while driving, cycling, or crossing
traffic. Avoid interactions under time pressure that demand attention. &
10 & 5 & 10 & 500 & HS3, HS7\\
\midrule
% ---------------------------------------------------------------- H2
\multicolumn{11}{l}{\textbf{H2: Lyric Display and Language Support}}\\
\midrule
H2.1 & Words in a script that is not Latin and has no spaces between words are segmented incorrectly. &
Learner studies something that is not a word or a wrong word boundary and may send it to Anki
(L2). &
a.~Tokenizer dictionary does not cover slang, fragments, or mixed scripts
found in lyrics.\newline
b.~Segmentation differs between editor and app versions. &
Comparison against a segmentation checked by the team for the acceptance suite
(PoC pass criterion R2). &
Compute segmentation in the editor, have the creator verify it line by line,
and store it in the song map. The learner app displays the stored segmentation
and never segments it again at runtime. &
6 & 8 & 4 & 192 & IN4\\
\midrule
H2.2 & Transliteration or reading is wrong (for example, a kanji with several
possible readings). &
Learner learns the wrong pronunciation (L2). &
a.~Automatic reading chooses the wrong reading for the context.\newline
b.~Creator enters a reading incorrectly. &
Creator verification step in the editor. Acceptance suite comparison. &
Store readings in the song map. Mark automatically generated readings as
\emph{unverified} until the creator confirms them, and block export of a map
with unverified lines. &
6 & 7 & 5 & 210 & IN4\\
\midrule
H2.3 & Lyrics in a supported script do not render correctly (missing glyphs or
wrong characters). &
Lyrics are unreadable or misleading (L2). &
a.~Device has no font for the script.\newline
b.~Text encoding is mishandled when the map is read. &
Rendering check of every acceptance suite line on both platforms. &
Bundle fonts that cover every script of the supported languages, and read and
write all song maps as UTF-8. &
5 & 3 & 2 & 30 & IN12\\
\midrule
H2.4 & Highlight changes flash rapidly across large areas of the screen. &
Discomfort or seizure risk for photosensitive learners (L1). &
a.~Fast songs with very short words.\newline
b.~Highlight style changes the background colour of whole lines. &
Review against WCAG~2.3.1 (three flashes or below threshold). Count highlight
changes per second in fast test songs. &
Use a highlight style that does not flash large areas (for example, an
underline or a moving marker). Limit large area changes to three per second,
and follow the operating system's reduce motion setting. &
9 & 3 & 4 & 108 & HS4\\
\midrule
H2.5 & Active line or word is shown only by colour, or text cannot be enlarged. &
Learners with low vision or colour vision deficiency cannot follow the lesson.
Strain from small text (L6). &
a.~Highlight relies on hue alone.\newline
b.~Fixed font sizes. &
Accessibility review with the OS text size setting at its largest and a
colour blindness simulator. &
Indicate the active line and word with a cue other than colour (weight, underline,
marker), and respect the device's text size setting. &
4 & 6 & 3 & 72 & HS5\\
\midrule
% ---------------------------------------------------------------- H3
\multicolumn{11}{l}{\textbf{H3: Annotation and Lesson Content}}\\
\midrule
H3.1 & A translation or contextual annotation is incorrect. &
Learner learns the wrong meaning and reinforces it through Anki reviews (L2). &
a.~Creator error or misunderstanding of slang or idiom.\newline
b.~Machine translation copied in without review. &
Review by a second team member before publishing. Learner error reports. &
Record the author and review status of each annotation. The editor does not
mark a map as published until every annotation has been reviewed by someone
other than its author. Let learners report a suspected error on any line. &
6 & 6 & 6 & 216 & IN5, AU1, AU2\\
\midrule
H3.2 & A lesson contains explicit, offensive, or culturally sensitive language
without warning. &
Learner is upset or uses an offensive term without knowing its register in
conversation (L2, L6). Content unsuitable for younger learners. &
a.~Authentic lyrics often contain profanity, slurs, or sensitive themes.\newline
b.~Annotations translate a term without describing how offensive it is. &
Creator checklist before publishing. Learner error reports. &
Each song map carries an explicit content flag. Flagged lessons show a warning
before opening and can be hidden through a learner setting. Annotations of
vulgar or offensive terms must include a usage note describing their register.
&
6 & 7 & 5 & 210 & HS6\\
\midrule
H3.3 & App design encourages compulsive use or discourages the learner. &
Stress, guilt, or loss of motivation (L6). &
a.~Streaks, loss penalties, or repeated notifications used to drive
engagement.\newline
b.~Learner opens lessons far above their level. &
Design review against HS7. Feedback from usability testing (Winter custom
document). &
Do not use streak loss penalties, countdowns, or unsolicited push
notifications. Show each lesson's difficulty level before it is opened. &
4 & 4 & 7 & 112 & HS7\\
\midrule
% ---------------------------------------------------------------- H4
\multicolumn{11}{l}{\textbf{H4: Song Map Format and Validator}}\\
\midrule
H4.1 & A map passes schema validation but its content is inconsistent. &
Desynchronization, missing annotations, or a crash during the lesson (L2, L3). &
a.~JSON Schema checks structure, not relationships (for example, overlapping or
decreasing timestamps, word times outside their line, times past the end of the
audio). &
Semantic validation in the editor on export and in the app on import.
Validator unit tests in CI. &
Add semantic rules on top of the schema: lines are in order and do not
overlap, word times fall inside their line, all times fall within the audio
duration, and every annotated target has both a translation and an
annotation (the full list is the SRS definition of a valid song map). &
6 & 5 & 3 & 90 & IN2\\
\midrule
H4.2 & Editor and app use different versions of the song map format. &
Fields are silently ignored or misread, or the app crashes (L2, L3). &
a.~Format changes during development.\newline
b.~Learner has an older app version. &
Version check on import. &
Each map declares its format version. The app refuses maps with an
unsupported major version and says why, instead of opening them partially. &
5 & 5 & 2 & 50 & IN3\\
\midrule
H4.3 & A crafted or malformed map is opened. &
App or editor freezes or crashes. Script injected through annotation text runs
in the web editor (L4). &
a.~Very large or deeply nested JSON.\newline
b.~Annotation text containing HTML or script.\newline
c.~Maps are files that can be shared and modified outside the editor. &
Fuzz and boundary tests of the parser. Size and depth checks before parsing.
&
Reject maps above a size and nesting depth limit before full parsing. Render
all map text as plain text in both the app and the editor. Never interpret it
as HTML. &
8 & 2 & 4 & 64 & IM1\\
\midrule
H4.4 & A map references audio or other resources by an external URL. &
Learner's device contacts an unknown server, revealing its IP address and
usage, or loads unexpected content (L4). &
a.~Creator links to hosted audio.\newline
b.~Map is modified by a third party. &
Validator rule on all resource paths. &
Allow only relative paths inside the map package. Reject absolute and remote
URLs. &
7 & 3 & 3 & 63 & IM2\\
\midrule
% ---------------------------------------------------------------- H5
\multicolumn{11}{l}{\textbf{H5: Song Map Editor}}\\
\midrule
H5.1 & Creator loses unsaved work. &
Lost timing and annotation work. The goal of creating a lesson in 45 minutes fails (L3).
&
a.~Browser tab closed or crashed.\newline
b.~Computer loses power. &
An indicator of unsaved changes in the editor. &
Autosave the draft to browser storage at least every 60 seconds and offer to
restore it when the editor reopens. Warn before the page closes with unsaved
changes. &
5 & 5 & 2 & 50 & IN6\\
\midrule
H5.2 & Creator accidentally overwrites or deletes timings or annotations. &
Corrupted map that may be exported without the creator noticing (L2, L3). &
a.~Bulk shift of timestamps applied to the wrong range.\newline
b.~Accidental delete. &
Preview playback before export. Export history (AU2). &
Provide undo and redo for all edits, and ask for confirmation before any bulk
operation that changes more than one line. &
5 & 5 & 5 & 125 & IN6, AU2\\
\midrule
H5.3 & Timestamps recorded in the editor are offset from the audio. &
Every lesson made on that computer is out of sync (L2). &
a.~Browser or Bluetooth audio output latency.\newline
b.~Creator reaction time when tapping along. &
Preview playback in the editor. Acceptance suite timing check. &
Provide a latency calibration step whose offset is applied to recorded
timestamps, and a preview mode that plays the map exactly as the app will. &
6 & 7 & 4 & 168 & IN7\\
\midrule
H5.4 & Copyrighted audio or lyrics are used without permission. &
Takedown requests or legal liability for the team and creators (L5). &
a.~Creator uses a commercial recording or lyrics found online.\newline
b.~Licence is not recorded, so it cannot be checked. &
Review of licence metadata before a map is published. &
Require source and licence fields for the audio and the lyrics before export,
and show the credits in the learner app. Limit the initial release to
content that is created by the team, in the public domain, or licensed, as in the Problem Statement.
&
9 & 4 & 5 & 180 & LG1\\
\midrule
% ---------------------------------------------------------------- H6
\multicolumn{11}{l}{\textbf{H6: Anki Handoff}}\\
\midrule
H6.1 & The same vocabulary item is added to Anki more than once. &
Cluttered deck and repeated reviews (L3). &
a.~Learner taps \emph{send} repeatedly.\newline
b.~Automatic retry after a timeout whose first request actually succeeded. &
Duplicate check before sending. Integration test that sends the same item
twice. &
Before adding a note, check the target deck for an existing note with the same
term and context. Disable the send control while a request is in progress. &
3 & 6 & 3 & 54 & IN9\\
\midrule
H6.2 & Note fields are filled incorrectly or text is garbled. &
Learner reviews a wrong or unreadable card (L2). &
a.~Fields assigned by position and the note type has a different order.\newline
b.~Text in a script other than Latin is not encoded as UTF-8. &
Check that the note type has the expected fields before sending. Acceptance
test of the \emph{Term}, \emph{Translation}, and \emph{Context} fields. &
Fill fields by name, not by position. Create the LingoLyrics note type if it is
missing and stop with a message if an existing one lacks the required fields.
Send all text as UTF-8. &
5 & 4 & 3 & 60 & IN10\\
\midrule
H6.3 & The app changes notes, decks, or note types other than its own. &
Learner's existing study data or review history is damaged (L3). &
a.~Programming error in the integration.\newline
b.~API calls with broader effects than intended. &
Code review of every Anki API call. Integration test against a test
collection. &
The app only adds notes, and only to the LingoLyrics deck. It never edits or
deletes existing notes, decks, or note types. &
8 & 2 & 5 & 80 & AC1\\
\midrule
H6.4 & Handoff fails silently. &
Learner believes the word was saved, and it is lost (L3). &
a.~Anki or AnkiDroid not installed or not running.\newline
b.~Permission denied.\newline
c.~App closed during the request. &
Response from Anki checked for every request. &
Report success only after Anki confirms the note was created. Keep failed
items in a local queue that survives an app restart, and allow retrying them or
exporting them to a file. &
6 & 5 & 6 & 180 & IN8\\
\midrule
H6.5 & AnkiConnect is configured to accept requests from any web page. &
Any website the learner visits can read or change their Anki collection (L3,
L4). &
a.~Setup instructions tell users to allow all origins (\texttt{*}) to make the
connection work. &
Review of setup documentation and default configuration. &
If AnkiConnect is chosen (PoC risk R3), allow only the specific origin the
product needs, never a wildcard, and use AnkiConnect's API key when it is
available. &
8 & 3 & 7 & 168 & AC2\\
\midrule
H6.6 & App requests platform permissions it does not need, or requests them
without explanation. &
Learner grants access they do not understand, or refuses and loses the feature
(L4). &
a.~Permissions declared and requested at install or first launch. &
Review of the app manifest and permission prompts. &
Request the AnkiDroid permission only when the learner first sends a word, with
a short explanation of why it is needed. Declare no other sensitive
permissions. &
4 & 4 & 2 & 32 & PR2\\
\midrule
% ---------------------------------------------------------------- H7
\multicolumn{11}{l}{\textbf{H7: Local Data Storage}}\\
\midrule
H7.1 & App crashes or is closed while vocabulary is waiting to be sent. &
Selected vocabulary is lost (L3). &
a.~Queue held only in memory.\newline
b.~Operating system terminates the app in the background. &
Restart test with items in the queue. &
Store the pending queue on the device as soon as an item is selected, and
remove an item only after Anki confirms it. &
6 & 4 & 5 & 120 & IN8\\
\midrule
H7.2 & Personal or usage data leaves the device. &
Privacy breach, including the learner's vocabulary history and listening
habits (L4). &
a.~Analytics or crash reporting libraries from other vendors.\newline
b.~Remote resources referenced by maps (H4.4). &
Review of dependencies and network traffic during testing. &
The initial release collects no personal data and sends no data off the device,
apart from notes the learner explicitly sends to Anki. Any analytics or crash
reporting added later must stay off unless the learner turns it on. &
8 & 2 & 6 & 96 & PR1, IM2\\
\end{longtable}
\end{landscape}
\ClearShipoutPictureFG

\begin{table}[h!]
\caption{High priority failure modes, in order of RPN} \label{TblPriority}
\begin{tabularx}{\textwidth}{lXrrrrl}
\toprule
\textbf{Ref.} & \textbf{Failure mode} & \textbf{S} & \textbf{O} & \textbf{D} &
\textbf{RPN} & \textbf{SR}\\
\midrule
H1.4 & Learner watches the screen while moving & 10 & 5 & 10 & 500 & HS3, HS7\\
H3.1 & Incorrect translation or annotation & 6 & 6 & 6 & 216 & IN5, AU1, AU2\\
H2.2 & Incorrect transliteration or reading & 6 & 7 & 5 & 210 & IN4\\
H3.2 & Explicit content without warning & 6 & 7 & 5 & 210 & HS6\\
H2.1 & Incorrect segmentation & 6 & 8 & 4 & 192 & IN4\\
H5.4 & Unlicensed audio or lyrics & 9 & 4 & 5 & 180 & LG1\\
H6.4 & Anki handoff fails silently & 6 & 5 & 6 & 180 & IN8\\
H5.3 & Editor timestamps offset from the audio & 6 & 7 & 4 & 168 & IN7\\
H6.5 & AnkiConnect accepts any web page & 8 & 3 & 7 & 168 & AC2\\
H1.2 & Audio starts unexpectedly or loudly & 8 & 5 & 4 & 160 & HS1, HS2\\
H2.4 & Rapid flashing & 9 & 3 & 4 & 108 & HS4\\
\bottomrule
\end{tabularx}
\end{table}

\section{Safety and Security Requirements} \label{SecRequirements}

The requirements below were discovered through the FMEA. Each requirement
lists the FMEA rows it mitigates. All of them are new requirements and have been
added to the SRS with the same identifiers, where each one also has a rationale
and a fit criterion: HS1 to HS7 under \emph{Safety-Critical Requirements}, AC, IN,
PR, AU, and IM under \emph{Security Requirements}, and LG1 under \emph{Legal
Requirements}. In the SRS, the numeric values used here (1 second, three
flashes per second, 60 seconds) are written as the symbolic constants
\texttt{DISCONNECT\_PAUSE}, \texttt{MAX\_FLASHES}, and
\texttt{AUTOSAVE\_INT}. The values of these and of \texttt{AUDIO\_DUR\_TOL},
\texttt{MAX\_MAP\_SIZE}, and \texttt{MAX\_NESTING} are listed in SRS
Appendix A, \emph{Symbolic Constants}. The SRS also lists, for each FMEA row, the requirements
that mitigate it (SRS Appendix B, \emph{Hazards to Requirements}).

\subsection{Health and Safety Requirements}

\begin{itemize}
  \item[HS1.] The learner app shall not start audio playback without an
        explicit action by the learner. \emph{(H1.2)}
  \item[HS2.] The learner app shall pause playback within 1 second when the
        active audio output device is disconnected. \emph{(H1.2)}
  \item[HS3.] The learner app shall provide a listening mode in which a
        lesson can be followed without looking at or touching the screen, and
        shall show a notice on first launch only, advising against using
        the screen while driving, cycling, or crossing traffic. \emph{(H1.4)}
  \item[HS4.] Lyric highlighting shall not produce more than three large area
        flashes in any period of one second, and shall follow the operating
        system's reduce motion setting. \emph{(H2.4)}
  \item[HS5.] The learner app shall indicate the active line and word with a
        cue other than colour alone, and shall respect the device's text size
        setting. \emph{(H2.5)}
  \item[HS6.] Each song map shall include an explicit content flag. Lessons
        flagged as explicit shall show a warning before opening and can be
        hidden through a learner setting. Annotations of vulgar or offensive
        terms shall include a note describing their register. \emph{(H3.2)}
  \item[HS7.] The learner app shall not use streak loss penalties, countdown
        timers, or unsolicited push notifications to encourage use, and shall
        show each lesson's difficulty level before it is opened.
        \emph{(H1.4, H3.3)}
\end{itemize}

\subsection{Access Requirements}

\begin{itemize}
  \item[AC1.] The learner app shall only add notes, and only to the
        LingoLyrics deck. It shall not edit or delete existing notes, decks,
        or note types in the learner's Anki collection. \emph{(H6.3)}
  \item[AC2.] If AnkiConnect is used, the product and its setup instructions
        shall allow only the specific origin the product requires, never a
        wildcard origin, and shall use an API key when AnkiConnect supports
        one. \emph{(H6.5)}
\end{itemize}

\subsection{Integrity Requirements}

\begin{itemize}
  \item[IN1.] The learner app shall drive lyric highlighting from the audio
        player's reported position and shall realign the highlight after
        every seek, pause, resume, and return from the background.
        \emph{(H1.1)}
  \item[IN2.] The song map validator shall check, in addition to the schema,
        the conditions of a valid song map given in the SRS
        (\emph{Formal Specification}, conditions V1 to V8): among them, lines
        are in order and do not overlap, word times fall within their line,
        all times fall within the audio duration, segmentation is verified,
        and every annotated target has both a translation and an annotation.
        The editor shall run it before export and the learner app shall run
        it on import.
        \emph{(H1.1, H4.1)}
  \item[IN3.] Each song map shall declare its format version. The learner app
        shall refuse to open a map with an unsupported major version and shall
        tell the learner why. \emph{(H4.2)}
  \item[IN4.] Song maps shall store word segmentation and transliteration
        explicitly, and the learner app shall display the stored values
        without segmenting again at runtime. The editor shall mark automatically
        generated values as unverified until the creator confirms them, and
        shall not export a map containing unverified lines.
        \emph{(H2.1, H2.2)}
  \item[IN5.] Each annotation shall record its author and review status. The
        editor shall not mark a map as published until every annotation has
        been reviewed by someone other than its author. \emph{(H3.1)}
  \item[IN6.] The editor shall autosave the current draft at least every 60
        seconds, offer to restore it when reopened, warn before closing with
        unsaved changes, provide undo and redo for all edits, and ask for
        confirmation before any operation that changes more than one line.
        \emph{(H5.1, H5.2)}
  \item[IN7.] The editor shall provide an audio latency calibration whose
        offset is applied to recorded timestamps, and a preview that plays the
        map as the learner app will. \emph{(H5.3)}
  \item[IN8.] The learner app shall report a vocabulary item as sent only
        after Anki confirms that the note was created. Items that have not
        been confirmed shall be stored on the device, kept across app
        restarts, and available to retry or export to a file.
        \emph{(H6.4, H7.1)}
  \item[IN9.] Before adding a note, the learner app shall check the
        LingoLyrics deck for an existing note with the same term and context,
        and shall disable the send control while a request is in progress.
        \emph{(H6.1)}
  \item[IN10.] The learner app shall fill Anki note fields by field name,
        shall stop with a message if the note type lacks the \emph{Term},
        \emph{Translation}, or \emph{Context} field, and shall send all text
        as UTF-8. \emph{(H6.2)}
  \item[IN11.] On import, the learner app shall check that the audio file
        referenced by a map exists, can be decoded, and has a duration within
        \texttt{AUDIO\_DUR\_TOL} of the duration recorded in the map, and
        shall reject the lesson with a message if not. \emph{(H1.3)}
  \item[IN12.] The learner app shall include fonts covering every script used
        by its supported languages, and all song maps shall be read and
        written as UTF-8. \emph{(H2.3)}
\end{itemize}

\subsection{Privacy Requirements}

\begin{itemize}
  \item[PR1.] The initial release shall not collect personal data or send any
        data off the device, other than notes the learner explicitly sends to
        Anki. Any analytics or crash reporting added later shall stay off unless the learner turns it on.
        \emph{(H7.2)}
  \item[PR2.] The learner app shall request the AnkiDroid permission only when
        the learner first sends a word, with an explanation of why it is
        needed, and shall request no other sensitive permissions.
        \emph{(H6.6)}
\end{itemize}

\subsection{Audit Requirements}

\begin{itemize}
  \item[AU1.] The learner app shall let a learner report a suspected error on
        any lyric line. The report shall contain the map identifier, map
        version, line, and the learner's comment, and shall be sent only when
        the learner chooses to send it (for example, through the project's
        issue tracker). \emph{(H3.1)}
  \item[AU2.] Each song map shall record a version number and, for every
        export, the date and the creator who exported it. \emph{(H3.1, H5.2)}
\end{itemize}

\subsection{Immunity Requirements}

\begin{itemize}
  \item[IM1.] The song map parser in both the editor and the learner app shall
        reject a map whose size exceeds \texttt{MAX\_MAP\_SIZE} or whose
        nesting depth exceeds \texttt{MAX\_NESTING} before parsing it fully,
        and both shall display all map text as plain text, never as HTML.
        \emph{(H4.3)}
  \item[IM2.] Song maps shall refer to audio and other resources only by
        relative paths inside the map package. The validator shall reject
        absolute and remote URLs. \emph{(H4.4, H7.2)}
\end{itemize}

\subsection{Legal Requirements}

\begin{itemize}
  \item[LG1.] The editor shall require source and licence information for the
        audio and lyrics before a map can be exported, and the learner app
        shall show these credits for each lesson. \emph{(H5.4)}
\end{itemize}

\section{Roadmap} \label{SecRoadmap}

Every high priority failure mode in Table~\ref{TblPriority} has at least one
mitigating requirement that will be implemented within the capstone timeline,
and all requirements needed for the goals in the Problem Statement will be
implemented within that timeline:

\begin{itemize}
  \item \textbf{Capstone timeline:} HS1, HS2, HS4, HS5, HS6, HS7, AC1, IN1 to IN4,
        IN6, IN8 to IN12, PR1, PR2, IM1, IM2, LG1, and AU2. HS7 and PR1 constrain
        the design rather than add features, so they are met by not adding
        engagement mechanics or analytics.
  \item \textbf{Depends on the Proof of Concept:} AC2 applies only if
        AnkiConnect is chosen for the Anki handoff (PoC risk R3). The detailed
        behaviour of IN8 to IN10 and PR2 will be confirmed once the handoff
        method is decided.
  \item \textbf{Partially implemented during the capstone:} IN5 will be met
        for the acceptance suite by a second team member reviewing each
        annotation, with the review status recorded in the map. Enforcing it
        in the editor is planned but may move to future work. IN7 will
        provide a manual offset. Automatic calibration is future work. HS3
        will provide the notice on first launch. A full listening mode
        depends on usability testing results.
  \item \textbf{Future work:} AU1 (error reporting in the app) requires a way to
        deliver reports, which the initial release, with no backend, does not
        have. Until then, reports are collected manually during usability
        testing.
\end{itemize}

The hazard analysis will be reviewed after the Proof of Concept and again
before the final demonstration to confirm which hazards have been mitigated
and which remain open.

\newpage{}

\section{Appendix:  Generative AI Use Disclosure}

\subsection{AI Tools Used}

\begin{itemize}
  \item Claude (Anthropic), model Claude Opus 5.5, through Claude Code.
\end{itemize}

\subsection{How and Where Used}

Claude was given the Problem Statement and Goals document, the Development
Plan, the course template and Hazard Analysis checklist, and three example
hazard analyses from previous capstone projects. It was used to:

\begin{itemize}

  \item Review the first draft of the document for grammatic errors and
        compare the document to previous student work and the checklist to 
        identify any missing sections or details.
  \item ensure the SRS was reviewed so the Hazard analysis did not miss any 
        relevant requirements 

\end{itemize}



\newpage{}

\section{Appendix:  Reflection}

\input{../Reflection.text}

\begin{enumerate}
    \item \textbf{What went well while writing this deliverable?}\\
    \textit{[To be completed by the team.]}

    \item \textbf{What pain points did you experience during this
    deliverable, and how did you resolve them?}\\
    \textit{[To be completed by the team.]}

    \item \textbf{Which of your listed risks had your team thought of before
    this deliverable, and which did you think of while doing this
    deliverable? For the latter ones (ones you thought of while doing the
    Hazard Analysis), how did they come about?}\\
    \textit{[To be completed by the team. The PoC risks R1 to R3 in the
    Development Plan correspond to H1.1, H2.1/H2.2, and H6. Most other rows are
    new to this deliverable.]}

    \item \textbf{Other than the risk of physical harm (some projects may not
    have any appreciable risks of this form), list at least 2 other types of
    risk in software products. Why are they important to consider?}\\
    \emph{Security and privacy risk:} software can expose or misuse personal
    data, or let an attacker run code or change data. For \progname{}, a
    shared song map is untrusted input and the app writes into the learner's
    Anki collection, so a flaw could damage data the learner has built up
    over years. These risks matter because users usually cannot see them and
    the harm is often permanent.

    \emph{Psychological and wellbeing risk:} engagement mechanics such as
    streaks and notifications can create guilt, stress, or compulsive use,
    and offensive content without warning can upset users. These risks
    matter because the people most affected, such as young or vulnerable
    users, are least able to recognize them, and they grow with the number
    of users.

    \emph{Legal and content accuracy risk:} distributing unlicensed content or
    teaching incorrect information can cause legal liability or lasting
    harm to users who trust the product. For a learning tool, incorrect
    content is the main way it can fail its users even when the software
    works perfectly.
\end{enumerate}

\end{document}
